\section{结论}
本文围绕多模态情感预测的时间组织、局部缺失与决策解释，建立三个相互关联的模型，主要结论如下。

第一，以原词为公共时间锚点，通过CTC强制对齐、子词归属与时间交叠池化，完成附件1全部100条样本的三模态特征提取。1932个原词中1920个生成时间区间，生成率99.38\%；全量结果表、观测掩码与典型片段给出特征与原始素材的对应关系，质量分量与边界扰动记录进一步刻画各锚点的可用程度，该时间表示为缺失掩码与证据回映提供了统一定义。

第二，采用EMT-DLFR教师与逆频率类别加权的三态感知学生模型，在验证集上取得0.6163的Macro-F1和0.5818的强度MAE，并完成附件3全部30条样本的预测。测试集基础Macro-F1为0.6189、MAE为0.6221；105种连续局部缺失条件下的平均Macro-F1为0.6066、MAE为0.6492。目标区间比例由5\%增至40\%时F1累计下降0.0240，缺口位置带来的平均差异为0.0090，缺口越重相对通用融合框架的优势越大，模型在局部缺失下保持稳定输出，性能损失集中在覆盖样本主体的文本长缺口上。

第三，模态子集训练与输入干预相结合的解释模型完成附件4全部20条样本的预测与解释，输出中性7条、负向7条、正向6条。冻结测试子集上分类与强度任务的高分片段删除优势分别为0.1318和0.2132，两种任务的局部排序均能定位影响当前输出的片段；118条可用证据完成时间与媒体回映，文本贡献总体占主导，少数非文本归因对解释基线较为敏感。

上述模型分别给出统一的时间表示、面向局部缺失的预测方式与对应模型决策的证据定位，形成从原始素材到预测与解释的可追溯链条。其结构不依赖特定数据集，可按新任务的特征维度、缺失模式与标签分布重新配置，用于具有转写文本的音视频情感分析场景。
